\section{Relation to existing work}
\label{sec:literature}\label{sec:literature-boundary}

The financial motivation connects costly information acquisition with the
maintenance of productive capabilities. \citet{VanNieuwerburghVeldkamp2010}
jointly model information acquisition and portfolio choice, showing how
learning incentives affect investment. \citet{EisfeldtPapanikolaou2013}
model organization capital embodied in key talent and link its firm-specific
productivity and outside options to financial risk. Our retention problem
focuses on declared reusable inputs to later strategy research: removing a
carrier can change which projects remain feasible even when its current
operating contribution is redundant. The closure is a specified capability
inventory, not an estimate of tacit knowledge or an institutional measure of
organization capital. This paper asks when a smaller library preserves that
inventory and its modeled uses; it does not determine whether every capability
is worth maintaining.

Belief-state reduction supplies the control architecture used here.
\citet{SmallwoodSondik1973} and \citet{Sondik1978} develop finite- and
infinite-horizon partially observed control; \citet{White1976} treats a
partially observed semi-Markov process. Our additional question concerns the
library component of state: which retained strategies leave both operating
choices and the declared generation process unchanged? The frontier--closure
pair summarizes that component under the factorization assumption. The belief
and active-project variables remain in the controlled state.

The preservation target distinguishes this problem from nearby reductions.
Incremental pruning removes redundant value vectors in a fixed POMDP backup
\citep{CassandraLittmanZhang1997}; state aggregation identifies behaviorally
equivalent states for a specified model \citep{GivanDeanGreig2003}. Neither
requires a removed strategy to remain available as a carrier for subsequent
generation. Compact control libraries trade static usefulness against library
size and search effort \citep{DeyEtAl2012}. Reward-free policy compression
uses set cover to represent state--action distributions for downstream tasks
\citep{MuttiDelColRestelli2022}. Here the choices are subsets of a given
source, and the preserved objects are its operating frontier and generative
closure.

Once identity closure makes those objects a finite tagged requirement set,
the algorithmic core is classical weighted cover. The complexity reduction
uses Karp's set-cover problem \citep{Karp1972}, and the greedy guarantee is
transferred from \citet{Chvatal1979}. The modeling contribution is the mapping
from the financial preservation condition to that problem; the generic cover
algorithms are established methods. Nonidentity closure requires a separate
representation of capability interactions. Maximum-coverage or adaptive
selection guarantees do not follow merely from calling capabilities a library.

In finance, automated construction has long used genetic search and
reinforcement methods \citep{AllenKarjalainen1999,MoodySaffell2001}.
Those studies concern constructing or operating strategies; our decision is
which entries in a maintained source can be removed while preserving its
specified opportunities. Our burden index measures the declared maintenance
requirements; it is not an estimate of institutional compliance costs.

A larger candidate menu also creates a larger selection problem. The
literature on data snooping and multiple testing explains why more available
strategies cannot itself establish better investment performance
\citep{SullivanTimmermannWhite1999,White2000,HarveyLiuZhu2016}. The financial
study therefore freezes a bootstrap-based proposal-time adoption rule before
held-out evaluation and reports abstention and adverse exercise. Its known
historical setting still limits external validity, and post-publication decay
provides a further reason for restraint \citep{McLeanPontiff2016}.
